% Einheit 4 von 5 – Vierecke nachweisen (bank/punkte-und-strecken-im-koordinatensystem/e4.jsonl, zusammenbau v0.1)
%% TODO Zeitmarke Einheit 4: Marken-Zeile nicht lesbar
%% TODO Zweigzeile Teil 1: was hier gelernt wird (Satz in Schülersprache aus dem Einheitstitel) – Einheit 4
\einheitenkopf[e4]{Einheit 4 von 5 · Vierecke nachweisen}
\zweigzeile{Abitur GK}
\setcounter{aufgabe}{16}

%% TODO Titel: Ich-kann-Satz fehlt in der Bank (Platzhalter: Kettenname) – Nr. 17
%% TODO Anweisung: ein Satz über der ersten Teilaufgabe fehlt in der Bank – Nr. 17
\begin{aufgabe}{Vierecke nachweisen}
%% TODO \rechenplatz{n}: Schrittzahl des Grundfalls fehlt in der Bank (nur bei fester Schreibform, 2.3 b)
\begin{teile}
\teil Nachzuweisen: $ABCD$ ist ein Rechteck. Was gehört zum Steckbrief? \\ \kreuz{vier gleich lange Seiten} \\ \kreuz{Parallelogramm und ein rechter Winkel} \\ \kreuz{ein Paar paralleler Seiten}
\teil Zeige: $A(1 | 2 | 0)$, $B(4 | 3 | 2)$, $C(3 | 6 | 3)$, $D(0 | 5 | 1)$ bilden ein Parallelogramm $ABCD$. \hfill (Abitur 2021 LK)
\teil Zeige: $A(2 | 1 | 0)$, $B(4 | 3 | 1)$, $C(2 | 4 | 3)$, $D(0 | 2 | 2)$ bilden ein Rechteck $ABCD$. \hfill (Abitur 2023 GK)
\teil Zeige: $A(0 | 1 | 3)$, $B(2 | 3 | 4)$, $C(3 | 5 | 6)$, $D(1 | 3 | 5)$ bilden eine Raute $ABCD$. \hfill (Abitur 2022 GK)
\teil Zeige: $A(2 | 2 | 2)$, $B(5 | 3 | 2)$, $C(6 | 5 | 4)$, $D(3 | 4 | 4)$ bilden ein Parallelogramm, aber kein Rechteck. \hfill (Abitur 2022 GK)
\teil Eine Kletterwand hat die Ecken $A(4 | 1 | 3)$, $B(1 | 4 | 3)$, $F(1 | 7 | 0)$ und $E(7 | 1 | 0)$ (1 LE = 1 m). Zeige: $ABFE$ ist ein Trapez, in dem zwei gegenüberliegende Seiten gleich lang sind. \hfill (Abitur 2018 GK)
\teil Eine Terrasse ist das Viereck $A(1 | -3 | 0)$, $B(7 | -1 | 0)$, $C(7 | 1 | 0)$, $D(1 | 3 | 0)$ (1 LE = 1 m). Zeige: $ABCD$ ist ein Trapez. Flächeninhalt? \hfill (Abitur 2026 GK) \\ \leerfeld[m²]
\teil Zeige: Das Viereck $K(0 | 0 | 2)$, $L(2 | 3 | 2)$, $M(6 | 0 | 2)$, $N(2 | -3 | 2)$ ist ein Drachenviereck. \hfill (Abitur 2023 GK)
\teil Eine ebene Rasenfläche hat die Ecken $A(2 | 1 | 0)$, $B(9 | 1 | 0)$, $C(8 | 4 | 0{,}5)$, $D(8 | 7 | 1)$ und $E(2 | 7 | 1)$ (1 LE = 1 m). Zeige: $AE$ ist parallel zu $CD$, und $CD$ und $DE$ bilden einen rechten Winkel. \hfill (Abitur 2021 GK)
\teil $P(2 | 1 | 0)$ und $Q(2 | 1 | 5)$ sind benachbarte Ecken eines Quadrats in der Ebene $3x_1 + 4x_2 = 10$. Gib eine weitere Ecke an. \hfill (Abitur 2022 GK)
\teil Ein Quadrat liegt in der $x_2x_3$-Ebene; $P(0 | 6 | 3)$ ist eine Ecke. Der Schnittpunkt seiner Diagonalen liegt auf der Geraden durch $(4 | 3 | 2)$ parallel zur $x_1$-Achse. Zeige: $Q(0 | 2 | 5)$ ist eine der beiden zu $P$ benachbarten Ecken. \hfill (Abitur 2018 LK)
\end{teile}
\end{aufgabe}

%% TODO Titel: Ich-kann-Satz fehlt in der Bank (Platzhalter: Kettenname) – Nr. 18
%% TODO Anweisung: ein Satz über der ersten Teilaufgabe fehlt in der Bank – Nr. 18
\begin{aufgabe}{Vierecke nachweisen – Fehler finden, Begründen}
\begin{teile}
\teil Ich finde den Fehler: $A(1 | 3 | 2)$, $B(4 | 4 | 4)$, $C(5 | 7 | 7)$, $D(2 | 6 | 5)$. Mara: „$\overrightarrow{AB} = (3 | 1 | 2)$, $\overrightarrow{CD} = (-3 | -1 | -2)$ – verschieden, also kein Parallelogramm.“
\teil Begründe: Ist $\overrightarrow{AB} \circ \overrightarrow{AD} \ne 0$, so ist das Parallelogramm $ABCD$ kein Rechteck – mehr muss man nicht prüfen.
\end{teile}
\end{aufgabe}
